% GENERATED FILE. Edit canonical lessons, then run npm run generate:books.
% canonical-source-hash: fnv1a64:64c493e0a087c328
% canonical-lessons: IT-C498-tata, IT-C498-pediatra, IT-C498-accudire, IT-C498-badare, IT-C498-ubbidire

\chapter{Chi bada ai bambini}
\label{ch:it-chi-bada-ai-bambini}
\hlchaptermodality{498}

\begin{chapteropening}
\textbf{By the end of this chapter:} \emph{I can talk about who looks after children.}

Complete the last lesson of chapter 498: I can talk about who looks after children.
\end{chapteropening}

\section[la tata]{la tata — nanny}

\label{lesson:IT-C498-tata}



\begin{quote}
Before the new one: say the Italian for to scold, then the Italian for punishment.
\end{quote}

\subsection*{You'll want to know: la tata}
\textbf{la tata} — \textquotedblleft{}nanny\textquotedblright{}.

\begin{quote}
La tata porta i bambini a scuola.
\end{quote}

\begin{quote}
\emph{The nanny takes the children to school.}
\end{quote}

Friendly speech; also \emph{la babysitter}.

\begin{cousinweb}
\emph{tata} is a word from child speech, made by doubling a syllable as small children do; the origin is not settled beyond that.
\end{cousinweb}

\begin{grammarlens}[title={What you've built}]
One more word for telling, explaining and describing at B1.
\end{grammarlens}

\subsection*{Your turn}
\begin{itemize}
\raggedright
  \item \emph{Say it:} \emph{la tata}
  \item \emph{Say it:} \emph{la tata}, once more
  \item \emph{Say it:} a sentence of your own with \emph{la tata}
  \item \emph{Recall:} say the Italian for to scold, then the Italian for punishment, then say \emph{la tata} again
\end{itemize}

\begin{tcolorbox}[breakable,colback=teal!4,colframe=teal!35!black,title={Before you move on}]
What does la tata mean? (\textquotedblleft{}Nanny\textquotedblright{}.) Say it once more.
\end{tcolorbox}
\section[il pediatra]{il pediatra — paediatrician}

\label{lesson:IT-C498-pediatra}



\begin{quote}
Before the new one: say the Italian for punishment, then the Italian for nanny.
\end{quote}

\subsection*{You'll want to know: il pediatra}
\textbf{il pediatra} — \textquotedblleft{}paediatrician\textquotedblright{}.

\begin{quote}
Domani porto mio figlio dal pediatra.
\end{quote}

\begin{quote}
\emph{Tomorrow I am taking my son to the paediatrician.}
\end{quote}

The same form for a woman: \emph{la pediatra}.

\begin{cousinweb}
\emph{pediatra} is a modern learned word built from Greek \emph{paîs, paidós} \textquotedblleft{}child\textquotedblright{} and \emph{iatrós} \textquotedblleft{}doctor\textquotedblright{} (EtymDB).
\end{cousinweb}

\begin{grammarlens}[title={What you've built}]
One more word for telling, explaining and describing at B1.
\end{grammarlens}

\subsection*{Your turn}
\begin{itemize}
\raggedright
  \item \emph{Say it:} \emph{il pediatra}
  \item \emph{Say it:} \emph{il pediatra}, once more
  \item \emph{Say it:} a sentence of your own with \emph{il pediatra}
  \item \emph{Recall:} say the Italian for punishment, then the Italian for nanny, then say \emph{il pediatra} again
\end{itemize}

\begin{tcolorbox}[breakable,colback=teal!4,colframe=teal!35!black,title={Before you move on}]
What does il pediatra mean? (\textquotedblleft{}Paediatrician\textquotedblright{}.) Say it once more.
\end{tcolorbox}
\section[accudire]{accudire — to care for, to look after}

\label{lesson:IT-C498-accudire}



\begin{quote}
Before the new one: say the Italian for nanny, then the Italian for paediatrician.
\end{quote}

\subsection*{You'll want to know: accudire}
\textbf{accudire} — \textquotedblleft{}to care for, to look after\textquotedblright{}.

\begin{quote}
Mia zia accudisce la nonna malata.
\end{quote}

\begin{quote}
\emph{My aunt looks after my sick grandmother.}
\end{quote}

An \emph{-isc-} verb: \emph{accudisco}, \emph{accudisci}.

\begin{cousinweb}
\emph{accudire} was borrowed, according to EtymDB, from Spanish \emph{acudir} \textquotedblleft{}to come to help\textquotedblright{}; the deeper origin is not settled.
\end{cousinweb}

\begin{grammarlens}[title={What you've built}]
One more word for telling, explaining and describing at B1.
\end{grammarlens}

\subsection*{Your turn}
\begin{itemize}
\raggedright
  \item \emph{Say it:} \emph{accudire}
  \item \emph{Say it:} \emph{accudire}, once more
  \item \emph{Say it:} a sentence of your own with \emph{accudire}
  \item \emph{Recall:} say the Italian for nanny, then the Italian for paediatrician, then say \emph{accudire} again
\end{itemize}

\begin{tcolorbox}[breakable,colback=teal!4,colframe=teal!35!black,title={Before you move on}]
What does accudire mean? (\textquotedblleft{}To care for, to look after\textquotedblright{}.) Say it once more.
\end{tcolorbox}
\section[badare]{badare — to look after; to mind}

\label{lesson:IT-C498-badare}



\begin{quote}
Before the new one: say the Italian for paediatrician, then the Italian for to care for.
\end{quote}

\subsection*{You'll want to know: badare}
\textbf{badare} — \textquotedblleft{}to look after; to mind\textquotedblright{}.

\begin{quote}
Chi bada ai bambini stasera?
\end{quote}

\begin{quote}
\emph{Who is looking after the children tonight?}
\end{quote}

\emph{Badare a} + person or thing; \emph{non badare a spese} — not to worry about the cost.

\begin{cousinweb}
\emph{badare} is inherited from Medieval Latin \emph{batāre} \textquotedblleft{}to gape, to stare\textquotedblright{} (EtymDB): first to stand and watch, then to look after.
\end{cousinweb}

\begin{grammarlens}[title={What you've built}]
One more word for telling, explaining and describing at B1.
\end{grammarlens}

\subsection*{Your turn}
\begin{itemize}
\raggedright
  \item \emph{Say it:} \emph{badare}
  \item \emph{Say it:} \emph{badare}, once more
  \item \emph{Say it:} a sentence of your own with \emph{badare}
  \item \emph{Recall:} say the Italian for paediatrician, then the Italian for to care for, then say \emph{badare} again
\end{itemize}

\begin{tcolorbox}[breakable,colback=teal!4,colframe=teal!35!black,title={Before you move on}]
What does badare mean? (\textquotedblleft{}To look after; to mind\textquotedblright{}.) Say it once more.
\end{tcolorbox}
\section[ubbidire]{ubbidire — to obey}

\label{lesson:IT-C498-ubbidire}



\begin{quote}
Before the new one: say the Italian for to care for, then the Italian for to look after.
\end{quote}

\subsection*{You'll want to know: ubbidire}
\textbf{ubbidire} — \textquotedblleft{}to obey\textquotedblright{}.

\begin{quote}
Il cane ubbidisce solo a mio padre.
\end{quote}

\begin{quote}
\emph{The dog only obeys my father.}
\end{quote}

An \emph{-isc-} verb with \emph{a}: \emph{ubbidire ai genitori}. Also \emph{obbedire}.

\begin{cousinweb}
\emph{ubbidire} is inherited from Latin \emph{oboedīre} \textquotedblleft{}to listen to, to obey\textquotedblright{} (EtymDB), built on \emph{audīre} \textquotedblleft{}to hear\textquotedblright{} (Lewis \& Short).
\end{cousinweb}

\begin{grammarlens}[title={What you've built}]
One more word for telling, explaining and describing at B1.
\end{grammarlens}

\subsection*{Your turn}
\begin{itemize}
\raggedright
  \item \emph{Say it:} \emph{ubbidire}
  \item \emph{Say it:} \emph{ubbidire}, once more
  \item \emph{Say it:} a sentence of your own with \emph{ubbidire}
  \item \emph{Recall:} say the Italian for to care for, then the Italian for to look after, then say \emph{ubbidire} again
\end{itemize}

\begin{tcolorbox}[breakable,colback=teal!4,colframe=teal!35!black,title={Before you move on}]
What does ubbidire mean? (\textquotedblleft{}To obey\textquotedblright{}.) Say it once more.
\end{tcolorbox}
